L'analyse de sentiment est une tâche fondamentale du traitement automatique du langage naturel (TALN).
Elle consiste à identifier la polarité émotionnelle exprimée dans un texte (positif, négatif, neutre).
Avec l'essor de Twitter, des millions de messages sont produits chaque jour, constituant une source
précieuse d'opinion publique.

Ce projet applique un pipeline complet de Text Mining sur des données Twitter réelles :

\begin{enumerate}[leftmargin=*,label=\arabic*.]
  \item \textbf{Collecte et compréhension des données} : analyse exploratoire du corpus Sentiment140.
  \item \textbf{Prétraitement NLP} : nettoyage, tokenisation, gestion des spécificités des tweets (emojis, hashtags).
  \item \textbf{Représentation vectorielle} : BoW, TF-IDF, N-grams, Word2Vec.
  \item \textbf{Modélisation} : comparaison de cinq modèles (Naïve Bayes, SVM, Random Forest, LSTM, BERT).
  \item \textbf{Évaluation et interprétation} : métriques, matrices de confusion, analyse d'erreurs.
\end{enumerate}

\begin{infobox}[Tâche de classification]
Étant donné un tweet $t$, prédire sa polarité $y \in \{0, 1\}$ où $0$ = sentiment négatif et $1$ = sentiment positif.
\end{infobox}
